%% sections/01-introduction.tex

\section{Introduction}
\label{sec:introduction}

Tabletop role-playing games (TRPGs) are collaborative narrative systems in which a
Dungeon Master (DM) designs and adjudicates a fictional world while players embody
characters who interact with that world~\cite{dungeons2014players}. The quality of a
TRPG session is a compound phenomenon: it depends simultaneously on mechanical
fairness (did the dice produce outcomes that felt earned?), narrative momentum (did
scenes move at the right pace?), character engagement (did each player-character
contribute meaningfully?), emotional landing (did designed payoffs register as
genuinely affecting?), and design fidelity (did the adventure deliver what it
promised?). Human participants navigate all of these dimensions implicitly during
play, and their retrospective judgments --- whether a session was ``good'' --- are
shaped by all of them at once~\cite{bowman2010functions}.

Measuring this compound quality systematically has long been recognized as difficult.
Most professional game development relies on qualitative exit surveys, post-session
interviews, and practitioner heuristics rather than structured, multi-dimensional
scoring instruments~\cite{playtesting2019, fullerton2008}. Academic treatments of
playtest methodology focus primarily on deficit detection --- the process of
identifying where the game breaks --- rather than on measuring the degree to which a
session succeeded across a defined quality space~\cite{zagal2008collaborative}.

The rise of large language model (LLM) systems capable of functioning as autonomous
DMs introduces a new version of this measurement problem that is both more tractable
and more urgent. When an AI system simultaneously plays the DM and all player
characters, the session is fully logged, dice rolls are reproducible, and every
narrative decision is documented. But human participants are absent: there are no
exit surveys to administer, no participant accounts to aggregate, and no experiential
reports to rely on. The only record of session quality is the session log itself ---
the structured narrative produced by the AI pipeline. Evaluating that log requires a
scoring instrument that can operate systematically on textual evidence alone.

This is the problem the Marathon Playtest Rubric was built to solve. The rubric is an
8-dimension scoring instrument, each dimension scored 0--10 via evidence-anchored
band descriptions, applied at the close of every session to produce a gate score out
of 80. It was developed and validated within the Marathon D\&D Workshop, a structured
research environment for designing and playtesting Dungeons \& Dragons 5e adventures
set in the Dragonlance campaign setting. All sessions were run by an AI system
(Claude, Anthropic) functioning as DM, with all player characters also simulated by
that system from detailed character-sheet specifications. The result is a corpus of
49 fully logged, fully reproducible AI-simulated sessions across 7 complete campaigns.

\subsection{The Measurement Problem for AI-Run TRPG}

Existing quality frameworks for TRPGs assume human participants and rely on
self-report~\cite{bowman2010functions, schut2013}. Player-type models such as
Bartle's taxonomy~\cite{bartle1996hearts} characterize player preferences but do not
provide session-level scoring. Immersion and affect frameworks~\cite{mekler2017framework,
calleja2011} require subjective reports of experiential states that simulated
player-characters cannot provide. Narrative coherence metrics from AI story generation
research~\cite{fan2019, ammanabrolu2021, riedl2016} evaluate whether generated text is
well-formed and plausible, but a session can be narratively coherent and still fail as
a TRPG experience: if player-characters never make meaningful choices, if designed
emotional payoffs do not land, or if dice mechanics and story systematically undermine
each other, the session has failed by TRPG standards regardless of how well-formed
the prose is.

Text evaluation metrics from natural language processing --- BLEU~\cite{papineni2002bleu},
ROUGE~\cite{lin2004rouge}, BERTScore~\cite{zhang2019bertscore} --- measure similarity
to a reference rather than the presence of specific quality properties. LLM-as-judge
approaches~\cite{zheng2023judging} are promising but introduce the circularity of
evaluating an AI system using the same class of system that produced the output.
Property-based evaluation --- checking that outputs contain evidence of specified
outcomes --- avoids this circularity and does not require a gold-standard
corpus~\cite{celikyilmaz2020}. The Marathon rubric adopts this approach: every score
requires specific textual evidence from the session log; general impressions are
explicitly prohibited by the scoring protocol.

\subsection{Why Session-Level Scoring Matters for TRPG Research}

Session-level scoring enables analyses that aggregate per-session observations do not.
By scoring each of 49 sessions on the same 8-dimensional instrument, we can track
quality trajectories across campaigns, compare opening sessions to closing sessions,
and compare campaigns with fundamentally different party archetypes on the same scale.
When the rubric is also an adaptive instrument --- amended when empirical use reveals
gaps --- the amendment history is itself a finding: it maps the quality space of
AI-simulated TRPG more precisely than any \textit{a priori} specification could.

Beyond TRPG specifically, the Marathon rubric offers a template for any research
context in which AI-generated creative outputs need systematic evaluation at scale
without participant feedback. The combination of property-based anchors, dimensional
independence, and forward-only amendment is applicable to AI-generated fiction,
AI-facilitated collaborative design, and AI-simulated social scenarios more broadly.

\subsection{Contributions}

This paper makes four contributions to the study of AI-simulated narrative game sessions:

\begin{enumerate}
  \item \textbf{The 8-dimension rubric specification.} We describe the complete rubric
  with per-dimension anchor language, band boundaries (0--3, 4--6, 7--9, 10), and
  scoring protocol, including cross-dimensional interaction anchors that emerged from
  empirical use. The rubric is described at sufficient detail to be adopted or adapted
  by other researchers.

  \item \textbf{A 49-session, 7-campaign validation corpus.} We report gate scores
  across all 49 sessions of the Marathon corpus, including per-dimension means, binding
  PASS rates, and the gate/panel spread metric that diagnoses structural artifacts in
  individual sessions. This is, to our knowledge, the largest published corpus of
  scored AI-simulated TRPG sessions.

  \item \textbf{An amendment-driven evolution trace.} We document all 13 rubric versions
  from v1.0 to v1.12, identifying the specific empirical events (session innovations)
  that triggered each amendment. This trace demonstrates how a property-based scoring
  rubric self-calibrates toward a maturing creative corpus.

  \item \textbf{Cross-campaign consistency evidence.} We show that gate scores are
  consistent across campaigns with fundamentally different archetypes (grief-themed,
  professional-code, duty-spine, party-of-strangers), confirming that the rubric
  measures session execution quality rather than thematic content.

  \item \textbf{Construct validity analysis.} We discuss construct validity through
  panel agreement analysis and scope remaining validity gaps requiring human
  inter-rater studies.
\end{enumerate}

\subsection{Paper Organization}

Section~\ref{sec:related} surveys related work in playtest methodology, AI as dungeon
master, quality measurement in creative and narrative systems, and expert review
simulation. Section~\ref{sec:methodology} describes the Marathon workshop pipeline in
full, the 8-dimension rubric with all anchor language, the amendment mechanism, and the
seven campaigns. Section~\ref{sec:evaluation} presents the empirical results: gate
score trajectories, per-dimension trends, version-controlled opening scores, cross-campaign
comparisons, and binding-threshold PASS rates. Section~\ref{sec:discussion} analyzes
what the 12-point improvement measures, discusses cross-campaign consistency, and
enumerates limitations and implications. Section~\ref{sec:conclusion} summarizes
contributions and identifies priorities for future work.
